% ============================================================
% doku_ki_kantenerkennung.tex
% Wissenschaftliche Dokumentation -- Methode D: KI-Kantenerkennung (HED)
% Bildbasierte Eisdickenmessung -- Fabian Rüth, TU Braunschweig
% Code: hed_eis.py
% ============================================================

\section{Methode D -- KI-Kantenerkennung (HED)}
\label{sec:hed}

Die vierte Methode ersetzt den handentworfenen Canny-Operator
(Kap.~\ref{sec:canny}) durch ein \textbf{gelerntes} Kantenmodell:
\emph{Holistically-Nested Edge Detection} (HED) \citep{xie2015hed}. HED ist ein
tiefes Faltungsnetz, das Kanten nicht aus einem festen Gradientenkriterium,
sondern aus antrainierten, mehrskaligen Merkmalen vorhersagt. Wie bei Methode~B
dient der Laser nur zum Croppen; Null- und Messlinie sind HED-Kanten.
Code: \texttt{hed\_eis.py}.

% ────────────────────────────────────────────────────────────
\subsection{Das HED-Netz}
\label{sec:hed_netz}

HED \citep{xie2015hed} basiert auf einem VGG-artigen Encoder mit fünf Stufen. An
jeder Stufe hängt ein \emph{Side-Output}, der eine Kantenkarte in der jeweiligen
Auflösung erzeugt; alle Side-Outputs werden auf die Bildgröße hochskaliert und zu
einer finalen Wahrscheinlichkeitskarte fusioniert. Durch die \emph{tiefe
Überwachung} (jede Stufe wird gegen die Kanten-Ground-Truth trainiert) erfasst
das Netz sowohl feine als auch grobe Kanten. Das Modell liefert pro Pixel eine
\textbf{Kantenwahrscheinlichkeit} $P_e \in [0,1]$, keine binäre Schwellentscheidung
wie Canny.

\paragraph{Einbindung.}
Das vortrainierte Caffe-Modell \citep{jia2014caffe} wird über das DNN-Modul von
OpenCV \citep{bradski2000} geladen (\texttt{readNetFromCaffe} mit
\texttt{deploy.prototxt} und \texttt{hed.caffemodel}). Zwei technische Details:
\begin{itemize}
  \item \textbf{Eigene \texttt{CropLayer}}: Die im Caffe-Netz verwendete
        \emph{Crop}-Schicht (zentriertes Zuschneiden der Side-Outputs auf die
        Zielgröße) ist in OpenCV nicht enthalten und wird als kleine
        Python-Schicht nachimplementiert und registriert.
  \item \textbf{Vorverarbeitung}: \texttt{blobFromImage} skaliert das Bild auf ein
        Vielfaches von~$16$ und subtrahiert die VGG-Trainingsmittelwerte
        $(104{,}0;\,116{,}7;\,122{,}7)$ pro Kanal -- notwendig, damit die
        Eingabestatistik der des Trainings entspricht.
\end{itemize}

% ────────────────────────────────────────────────────────────
\subsection{Kantenkarte und Referenzlinie}
\label{sec:hed_karte}

Die HED-Ausgabe (Funktion \texttt{hed\_karte}) wird auf die Crop-Größe skaliert,
mit $255$ multipliziert und bei $T_\mathrm{HED} = 40$ geschwellt, auf das ROI-Band
beschränkt und anschließend \textbf{skelettiert} \citep{zhang1984}, da HED-Kanten
mehrere Pixel breit sind:
\begin{equation}
  E = \operatorname{skeleton}\bigl( (255\,P_e > T_\mathrm{HED}) \wedge \text{Band}\bigr).
\end{equation}
Aus $E$ von Frame~0 wird -- mit derselben Verkettungs-Infrastruktur wie in
Kap.~\ref{sec:canny_referenz} (\texttt{skelett\_teilstuecke},
\texttt{fit\_teilstueck}, \texttt{verbinde}, \texttt{bogenlaenge\_und\_normalen})
-- die geordnete \textbf{Nulllinie} $d_0(s)$ mit Bogenlänge und Außennormalen
gebildet. HED arbeitet direkt auf den bereits gecroppten PNGs, ein erneutes Laden
der großen TIFF-Dateien entfällt.

% ────────────────────────────────────────────────────────────
\subsection{Detektion mit Sprungbegrenzung}
\label{sec:hed_detektion}

Wie bei Canny (Kap.~\ref{sec:canny_tracking}) wird je Station die nächste
Kantenkreuzung zur Vorframe-Lage gewählt. Zusätzlich begrenzt eine
\textbf{Sprungschwelle} die zulässige Änderung pro Frame:
\begin{equation}
  d_N(s_i) =
  \begin{cases}
    k^\ast, & |k^\ast - d^{\text{prev}}(s_i)| \le \Delta_{\max}=5\,\text{px},\\
    \text{NaN}, & \text{sonst},
  \end{cases}
  \qquad
  k^\ast = \arg\min_{k_g} |k_g - d^{\text{prev}}(s_i)| .
\end{equation}
Ohne diese Begrenzung springt die Detektion in den ersten Frames auf die
$\approx 50\,\text{px}$ weiter außen liegende Kante des hell beleuchteten Bandes
und meldet Scheineis; $\Delta_{\max}$ erzwingt den physikalisch plausiblen,
langsamen Eiszuwachs. Eisdicke, Median-Glättung und Kalibrierung folgen wie in
Gl.~\eqref{eq:dicke}.

% ────────────────────────────────────────────────────────────
\subsection{Ergebnisse, Laufzeit und Grenzen}
\label{sec:hed_ergebnisse}

HED liefert den Endwert $\approx 5{,}5\,\text{px} \approx 0{,}40\,\text{mm}$ und
erzeugt die visuell sauberste, durchgehendste Kante auf der Laser-/Eislinie -- die
Stärke des gelernten Modells gegenüber Canny.

\paragraph{Laufzeit.}
Der Vorwärtsdurchlauf kostet auf der CPU $\approx 3{,}7\,\text{s}$ pro
Vollbild-Frame ($\approx 0{,}3$\,FPS). HED ist damit zwar die genaueste
Kantenmethode, für eine \textbf{Live-Messung} aber ungeeignet; hier sind
Methode~A/B ($\approx 11$\,FPS) oder das U-Net (Kap.~\ref{sec:semseg},
$\approx 75$\,FPS) vorzuziehen. Auf einer GPU wäre HED um etwa eine
Größenordnung schneller.

\paragraph{Einordnung.}
Mit HED liegen vier methodisch unabhängige Verfahren vor -- klassischer
Intensitätsrücken (A), klassische Kante (B), gelernte Segmentierung (C) und
gelernte Kante (D) --, deren Wachstumskurven auf $\pm 0{,}5\,\text{px}
\approx \pm 0{,}04\,\text{mm}$ übereinstimmen. Der Ensemble-Median beträgt
$\approx 6{,}2\,\text{px} \approx 0{,}45\,\text{mm}$; die Übereinstimmung
unabhängiger Prinzipien ist die zentrale Validierung der Messkette.
